\chapter*{Mở đầu}
\addcontentsline{toc}{chapter}{\vspace*{-8pt}  Mở đầu} 

Trong lĩnh vực khoa học máy tính, lập trình hợp ngữ Assembly là một trong những kỹ năng nền tảng quan trọng, giúp người học hiểu sâu về cách thức hoạt động của máy tính ở mức thấp. Ngôn ngữ Assembly là ngôn ngữ lập trình cấp thấp, gần với mã máy nhưng vẫn có thể đọc hiểu được đối với con người. Việc nắm vững Assembly giúp lập trình viên hiểu rõ hơn về kiến trúc máy tính, quản lý bộ nhớ, và tối ưu hóa hiệu suất chương trình.

Báo cáo này tập trung vào hai phần chính: phần thứ nhất trình bày các bài tập lập trình hợp ngữ Assembly, và phần thứ hai phân tích khảo sát hệ thống bộ nhớ. Trong phần lập trình Assembly, chúng tôi thực hiện nhiều bài tập đa dạng như chuyển đổi ký tự thường thành hoa, chuyển đổi giữa các hệ cơ số, tính toán số học, và xử lý chuỗi. Mỗi bài tập đều được minh họa bằng mã nguồn và kết quả thực thi, giúp người đọc hiểu rõ cách thức hoạt động của từng chương trình.

Phần phân tích khảo sát hệ thống bộ nhớ tập trung vào việc tìm hiểu cấu hình máy tính và cơ chế quản lý bộ nhớ. Chúng tôi sử dụng các công cụ như CPU-Z để khảo sát thông số kỹ thuật của CPU, cache, RAM, và các thiết bị lưu trữ. Đồng thời, chúng tôi cũng sử dụng công cụ debug để theo dõi và phân tích nội dung các thanh ghi như IP, DS, ES, SS, CS, BP, SP, từ đó hiểu rõ hơn về cơ chế quản lý bộ nhớ của hệ thống.

Mục tiêu của báo cáo này là giúp người đọc hiểu sâu hơn về lập trình Assembly và kiến trúc máy tính, đồng thời cung cấp các ví dụ thực tế về cách thức áp dụng kiến thức này vào việc phát triển và tối ưu hóa phần mềm. Chúng tôi hy vọng rằng báo cáo này sẽ là tài liệu tham khảo hữu ích cho sinh viên và những người quan tâm đến lĩnh vực lập trình cấp thấp và kiến trúc máy tính.
